\luaexec{dofile('code/protocol.lua')}

\begin{table}[H]
    \centering
    \caption{Результаты измерений частотомером ЦЧ}
    \label{tab:freq_counter_raw}
    \begin{tblr}[expand=\directlua]{
        colspec = { Q[c,m] Q[c,m] X[c,m] },
        hlines, vlines,
    }
        {№ вар.} & {Время счета $t_{\text{сч}}$, с} & {Показания $f_x$, Гц} \\
        \SetCell[r=3]{m} 1 & \val*[1]{p.fc.o1.t[1]} & \val*[3]{p.fc.o1.f_x[1]} \\
                           & \val*[1]{p.fc.o1.t[2]} & \val*[3]{p.fc.o1.f_x[2]} \\
                           & \val*[1]{p.fc.o1.t[3]} & \val*[3]{p.fc.o1.f_x[3]} \\
        \midrule
        \SetCell[r=3]{m} 2 & \val*[1]{p.fc.o2.t[1]} & \val*[3]{p.fc.o2.f_x[1]} \\
                           & \val*[1]{p.fc.o2.t[2]} & \val*[3]{p.fc.o2.f_x[2]} \\
                           & \val*[1]{p.fc.o2.t[3]} & \val*[3]{p.fc.o2.f_x[3]} \\
        \midrule
        \SetCell[r=3]{m} 3 & \val*[1]{p.fc.o3.t[1]} & \val*[3]{p.fc.o3.f_x[1]} \\
                           & \val*[1]{p.fc.o3.t[2]} & \val*[3]{p.fc.o3.f_x[2]} \\
                           & \val*[1]{p.fc.o3.t[3]} & \val*[3]{p.fc.o3.f_x[3]} \\
        \midrule
        \SetCell[r=3]{m} 4 & \val*[1]{p.fc.o4.t[1]} & \val*[3]{p.fc.o4.f_x[1]} \\
                           & \val*[1]{p.fc.o4.t[2]} & \val*[3]{p.fc.o4.f_x[2]} \\
                           & \val*[1]{p.fc.o4.t[3]} & \val*[3]{p.fc.o4.f_x[3]} \\
    \end{tblr}
\end{table}

\begin{table}[H]
    \centering
    \caption{Результаты измерений параметров сигнала осциллографом}
    \label{tab:osc_freq_raw}
    \begin{tblr}[expand=\directlua]{
        colspec = { Q[c,m] X[c,m] X[c,m] },
        hlines, vlines,
    }
        {№ изм.} & {Коэф. развертки $k_p$, мс/дел} & {Размер изображения $L_T$, дел} \\
        1 & \val*[1m-3]{p.of.td[1]} & \val*[1]{p.of.L_T[1]} \\
        2 & \val*[1m-3]{p.of.td[2]} & \val*[1]{p.of.L_T[2]} \\
        3 & \val*[1m-3]{p.of.td[3]} & \val*[1]{p.of.L_T[3]} \\
        4 & \val*[1m-3]{p.of.td[4]} & \val*[1]{p.of.L_T[4]} \\
    \end{tblr}
\end{table}

\begin{table}[H]
    \centering
    \caption{Результаты измерений фазового сдвига двумя методами}
    \label{tab:phase_raw}
    \begin{tblr}[expand=\directlua]{
        colspec = { Q[c,m] X[c,m] X[c,m] X[c,m] X[c,m] X[c,m] X[c,m] },
        hlines, vlines,
    }
        \SetCell[r=2]{m} {№} & \SetCell[c=2]{c} {Лин. развертка} &                       & \SetCell[c=2]{c} {Лиссажу} &                     & \SetCell[r=2]{m} {Развертка, мс/дел} & \SetCell[r=2]{m} {Отклонение, В/дел} \\
                             & {$\tau$, дел}                     & {$L_T$, дел}          & {$A$, дел}                 & {$B$, дел}          &                                       &                                       \\
        1                    & \val*[1]{p.ph.tau[1]}             & \val*[1]{p.ph.L_T[1]} & \val*[1]{p.ph.A[1]}        & \val*[1]{p.ph.B[1]} & \val*[1m-3]{p.ph.td[1]}               & \val*[1]{p.ph.vd[1]}                  \\
        2                    & \val*[1]{p.ph.tau[2]}             & \val*[1]{p.ph.L_T[2]} & \val*[1]{p.ph.A[2]}        & \val*[1]{p.ph.B[2]} & \val*[1m-3]{p.ph.td[2]}               & \val*[1]{p.ph.vd[2]}                  \\
        3                    & \val*[1]{p.ph.tau[3]}             & \val*[1]{p.ph.L_T[3]} & \val*[1]{p.ph.A[3]}        & \val*[1]{p.ph.B[3]} & \val*[1m-3]{p.ph.td[3]}               & \val*[1]{p.ph.vd[3]}                  \\
        4                    & \val*[1]{p.ph.tau[4]}             & \val*[1]{p.ph.L_T[4]} & \val*[1]{p.ph.A[4]}        & \val*[1]{p.ph.B[4]} & \val*[1m-3]{p.ph.td[4]}               & \val*[1]{p.ph.vd[4]}                  \\
    \end{tblr}
\end{table}

Заданная преподавателем частота $f = \val[0]{p.f}$ Гц.